\chapter[4.2.a--b Diseño del control inteligente]{4.2.a--b Diseño del sistema de control inteligente}
\label{ch:diseno}

\section{4.2.a Selección y justificación de la estrategia de control}

El sistema R2ET reúne tres rasgos incómodos para un control clásico: una no linealidad leve (la saturación del actuador), perturbaciones de amplitud conocida pero duración variable, y dos variables de proceso acopladas por el reparto del robot. Nada de esto impide usar un PID. Lo que hace es limitar su techo, porque un PID no coordina esteras por sí mismo ni reajusta sus parámetros cuando cambia el peso. De ahí que, junto al PID de referencia, se evalúen tres alternativas inteligentes.

\begin{longtable}{@{}P{0.12\linewidth}P{0.22\linewidth}P{0.55\linewidth}@{}}
\caption{Justificación de cada método de control para el sistema R2ET.}
\label{tab:justificacion}\\
\toprule
Método & Tipo & Justificación para R2ET \\ \midrule \endfirsthead
\toprule Método & Tipo & Justificación \\ \midrule \endhead
PID & Clásico: referencia & Elimina el error estacionario por acción integral. Bien conocido, fácil de auditar y mantener. Limitación: un PID escalar por estera no coordina el reparto del robot ni adapta parámetros ante la perturbación de peso.\\[4pt]
Fuzzy & Reglas lingüísticas & Las reglas Mamdani son auditables sin conocimiento de sistemas de control avanzados. No necesita datos históricos de la planta. Incorpora directamente el conocimiento de un experto sobre cuándo acelerar, frenar o redirigir el robot \cite{ross2010fuzzy,kasabov1996foundations}.\\[4pt]
Deep Learning & Optimización directa sobre planta & Entrena la política minimizando el IAE directamente sobre la planta mediante el método adjunto (co-estado). No requiere controlador externo de referencia. La inferencia en despliegue es de coste computacional fijo por ciclo \cite{reinoso2022redes}.\\[4pt]
Q-Learning & Refuerzo tabular & Aprende la política por interacción con el entorno, sin modelo de la planta. Puede continuar actualizando la tabla durante la operación (aprendizaje en línea). Es el método más adecuado cuando la eficiencia energética es el criterio dominante \cite{sutton2018reinforcement}.\\
\bottomrule
\end{longtable}

Cada método gana en un terreno distinto. El difuso es el más indicado para un despliegue auditable, porque sus 35~reglas las lee y las retoca el personal de planta sin herramientas de aprendizaje automático. El PID retuneado, en cambio, marca el menor IAE de la coordinación de dos esteras. El Q-Learning se lleva la eficiencia energética y la robustez. Y la red profunda combina lo mejor de ambos mundos: el mejor compromiso del lazo servo (IAE la mitad que el PID en Fase~1, con el mando más suave y sin saturar) y, en la coordinación, el mayor margen de seguridad y productividad.

\section{4.2.b Variables, universos de discurso y señales de retroalimentación}

\subsection*{Variables lingüísticas}

\begin{longtable}{@{}P{0.20\linewidth}P{0.18\linewidth}P{0.18\linewidth}P{0.32\linewidth}@{}}
\caption{Variables lingüísticas, universos de discurso y etiquetas del controlador.}
\label{tab:variables_ling}\\
\toprule
Variable & Universo & Etiquetas & Uso \\ \midrule \endfirsthead
\toprule Variable & Universo & Etiquetas & Uso \\ \midrule \endhead
\multicolumn{4}{@{}l}{\textit{Entradas}} \\ \midrule
$e = r - y$ (servo) & $[-1.0,\;1.0]$ & NL, NM, NS, ZE, PS, PM, PL & Error de seguimiento en Fase~1\\
$e_i = n_i - n^*$ & $[-1.5,\;1.5]$~piezas & NL, NM, NS, ZE, PS, PM, PL & Error de ocupación de estera $i$\\
$\Delta e$ & $[-0.5,\;0.5]$ & NB, NS, ZE, PS, PB & Tendencia del error\\
$d$ & $[0,\;1]$ & SIN, CON & Cambio de peso activo\\
$b = n_1 - n_2$ & $[-2.0,\;2.0]$~piezas & E1-pesada, Balanceada, E2-pesada & Desequilibrio entre esteras\\
\midrule
\multicolumn{4}{@{}l}{\textit{Salidas}} \\ \midrule
$u_i$ & $[0.05,\;0.98]$~(norm.) & 7 niveles (FAM por nivel de error/derivada) & Velocidad de descarga de estera $i$\\
$r_1$ & $[0.20,\;0.80]$\newline(fracción) & Bajo, Medio, Alto & Fracción del robot hacia E1\\
$a$ & $[0.72,\;1.00]$\newline(fracción) & — & Admisión (gestionada por supervisor)\\
\bottomrule
\end{longtable}

\subsection*{Señales de retroalimentación}

Fase~1: salida $y(k)$, error $e(k)=r-y(k)$, derivada $\Delta e(k)=e(k)-e(k-1)$, y perturbación $d(k)$ medida (o estimada vía observador de Luenberger en implementación real).

Fase~2: niveles $n_1(k)$ y $n_2(k)$ medidos por sensores de presencia de piezas, errores de ocupación $e_i(k)=n_i(k)-n^*$, variaciones $\Delta n_i(k)=n_i(k)-n_i(k-1)$, desequilibrio $b(k)=n_1-n_2$, y señal $d(k)$ de atasco.

\section{PID discreto: línea base}

\[
u(k) = u_{ff} + K_p\,e(k) + K_i\,\sigma(k) + K_d\,\Delta e(k)
\]

La integral $\sigma(k)$ usa anti-windup condicional: sólo se actualiza cuando el actuador no está saturado en la misma dirección que el error:
\[
\sigma(k+1) = \operatorname{sat}_{[-\sigma_{max},\,\sigma_{max}]}\!\bigl(\sigma(k) + e(k)\bigr)
\quad\text{si }u(k)\text{ no satura en sentido de }e
\]

Fase~1: $K_p=4.0$, $K_i=0.30$, $K_d=1.0$, $u_{ff}=0.60$, $\sigma_{max}=6.0$. Las ganancias son altas porque, al ser $b=0.08$ (planta de ganancia DC unitaria), el lazo necesita más acción proporcional para igual autoridad de control; se ajustaron por búsqueda minimizando el IAE con sobreimpulso nulo.

Fase~2: $K_p=1.3$, $K_i=0.18$, $K_d=0.28$, $u_{ff}=u_{ss}=0.457$. El supervisor ajusta $K_p$ y $K_i$ dinámicamente según el nivel de congestión (Capítulo~\ref{ch:supervisor}).

\section{Controlador difuso Mamdani}

\subsection*{Base de conocimiento}

La base de conocimiento del controlador difuso tiene dos componentes \cite[cap.~5]{passino2005biomimicry}:
\begin{itemize}
  \item Base de datos: las variables lingüísticas, sus universos de discurso y las funciones de pertenencia (Tabla~\ref{tab:variables_ling} y la estrategia de fuzzificación siguiente).
  \item Base de reglas: las 35 reglas $(e,\Delta e)\!\to\!u$ más la regla de perturbación, codificadas como matriz asociativa difusa (FAM, Tabla~\ref{tab:reglas_fuzzy}).
\end{itemize}
Sobre esta base, el motor de inferencia aplica la conjunción AND~=~mínimo y la defuzzificación convierte el grado de activación de las reglas en la velocidad de mando. Los tres pasos (fuzzificación, inferencia, defuzzificación) se detallan a continuación.

\subsection*{Estrategia de fuzzificación: funciones de membresía}

\artifactfigure{figures/plots/fig_fuzzy_mf.png}{0.96}{Funciones de pertenencia: error $e$ (7~conjuntos triangulares/trapezoidales), derivada $\Delta e$ (5~conjuntos), perturbación $d$ (2~conjuntos) y velocidad de salida por nivel de error para Fase~1 y Fase~2.}

Se emplean 7 conjuntos para el error (resolución densa cerca de cero) y 5 para la derivada, lo que da una base de 35 reglas (frente a las 15 de una partición $5\times3$) y un control más fino cerca del punto de operación. El diseño de las funciones de pertenencia, la inferencia y la defuzzificación sigue el procedimiento estándar de control basado en reglas \cite[cap.~5]{passino2005biomimicry}, \cite{ross2010fuzzy}.

Nomenclatura de las etiquetas. Cada etiqueta lingüística combina el signo del valor con su magnitud, según la convención habitual en control difuso: el prefijo es el signo (N~=~negativo, P~=~positivo) y el sufijo la magnitud (S~=~pequeña, \emph{small}; M~=~media, \emph{medium}; L/B~=~grande, \emph{large}/\emph{big}); ZE es el conjunto centrado en cero. Así, para el error: NL/NM/NS son negativo grande/medio/pequeño, ZE es cero y PS/PM/PL son positivo pequeño/medio/grande; para la derivada, NB y PB son negativo y positivo grandes. El significado físico es directo: un error PL (estera muy por encima del objetivo $n^*$) pide mayor velocidad de descarga, mientras que un error NL (estera casi vacía) pide frenarla; la derivada $\Delta e$ matiza esa acción según la tendencia del nivel.

Para el error $e \in [-1.5,\;1.5]$:
\begin{align*}
\mu_{NL} &= \text{trap}(-1.5,-1.5,-0.7,-0.4) & \mu_{NM} &= \text{tri}(-0.6,-0.35,-0.12) & \mu_{NS} &= \text{tri}(-0.22,-0.11,0) \\
\mu_{ZE} &= \text{tri}(-0.07,0,0.07) & \mu_{PS} &= \text{tri}(0,0.11,0.22) & \mu_{PM} &= \text{tri}(0.12,0.35,0.6) \\
\mu_{PL} &= \text{trap}(0.4,0.7,1.5,1.5) & &
\end{align*}

Para la derivada $\Delta e \in [-0.5,\;0.5]$:
\[
\mu_{NB}=\text{trap}(-0.5,-0.5,-0.22,-0.10),\;\; \mu_{NS}=\text{tri}(-0.18,-0.08,0),\;\; \mu_{ZE}=\text{tri}(-0.05,0,0.05),
\]
\[
\mu_{PS}=\text{tri}(0,0.08,0.18),\;\; \mu_{PB}=\text{trap}(0.10,0.22,0.5,0.5)
\]

\subsection*{Reglas de inferencia (base de reglas: 35 + 1 de perturbación)}

El consecuente de cada regla se sintetiza por niveles lingüísticos, equivalente a una matriz asociativa difusa (FAM):
\[
c(e,\Delta e) = \operatorname{sat}\!\big(u_{base} + L_e\cdot 0.12 + L_{\Delta e}\cdot 0.05,\;0.05,\;0.98\big),
\]
con niveles de error $L_e\in\{-3,\ldots,3\}$ (NL$\to$PL) y de derivada $L_{\Delta e}\in\{-2,\ldots,2\}$ (NB$\to$PB). La Tabla~\ref{tab:reglas_fuzzy} muestra la velocidad resultante para Fase~1 ($u_{base}=0.6$).

\begin{longtable}{@{}P{0.10\linewidth}P{0.08\linewidth}P{0.08\linewidth}P{0.08\linewidth}P{0.08\linewidth}P{0.08\linewidth}@{}}
\caption{FAM de 35 reglas: velocidad de salida $c$ para Fase~1 ($u_{base}=0.6$) según $(e,\Delta e)$.}
\label{tab:reglas_fuzzy}\\
\toprule
$e \backslash \Delta e$ & NB & NS & ZE & PS & PB \\ \midrule \endfirsthead
\toprule $e \backslash \Delta e$ & NB & NS & ZE & PS & PB \\ \midrule \endhead
NL & 0.14 & 0.19 & 0.24 & 0.29 & 0.34 \\
NM & 0.26 & 0.31 & 0.36 & 0.41 & 0.46 \\
NS & 0.38 & 0.43 & 0.48 & 0.53 & 0.58 \\
ZE & 0.50 & 0.55 & 0.60 & 0.65 & 0.70 \\
PS & 0.62 & 0.67 & 0.72 & 0.77 & 0.82 \\
PM & 0.74 & 0.79 & 0.84 & 0.89 & 0.94 \\
PL & 0.86 & 0.91 & 0.96 & 0.98 & 0.98 \\
\bottomrule
\end{longtable}

Regla de perturbación: SI $d=\text{CON}$ ENTONCES $u=u_{base}+0.30$, con intensidad $\mu_{CON}(d)$, activa sólo cuando $d>0.35$ (umbral de la MF trapecial). Se combina en paralelo con las 35 reglas en el promedio ponderado, reforzando la descarga ante el atasco sin interferir con el régimen nominal.

La inferencia usa AND = mínimo (intersección de Zadeh): el grado de activación de cada regla queda limitado por la premisa más débil \cite[cap.~5]{passino2005biomimicry}.

\subsection*{Estrategia de defuzzificación}

La salida es el promedio ponderado de los consecuentes (WA), método de defuzzificación estándar para singletons \cite[cap.~5]{passino2005biomimicry}:
\[
u(k) = \operatorname{sat}\!\left(\frac{\sum_j w_j\,c_j}{\sum_j w_j + \varepsilon},\;0.05,\;0.98\right),
\]
donde $w_j$ es la activación de cada regla y $\varepsilon=10^{-9}$ evita la división por cero cuando ninguna regla activa (en ese caso la salida se inicializa a $u_{base}$). Para Fase~2 se usa $u_{base}=0.457$, lo que desplaza toda la tabla al punto de operación de las esteras.

\section{Red neuronal profunda (Deep Learning)}

La red es un perceptrón multicapa que actúa como aproximador no lineal de la política de control \cite[cap.~4]{passino2005biomimicry}. Ambas redes se entrenan por optimización directa de política sobre la planta: la red minimiza el IAE directamente sobre la dinámica del sistema sin ningún controlador externo de referencia. El mecanismo de entrenamiento es el método adjunto (co-estado discreto), que calcula el gradiente del IAE respecto a los pesos de la red propagando el error a través de la dinámica de la planta hacia atrás en el tiempo.

\subsection*{Método adjunto (fundamento teórico)}

Dado el sistema $y(k+1)=A\,y(k)+B\,u(k)$ y la pérdida IAE $L=\sum|y(k)-r|$, el co-estado $\lambda(k)$ se calcula hacia atrás desde $T$:

\begin{align}
\lambda(T) &= \operatorname{sign}\!\left[y(T)-r\right]\,/\,T, \quad
\lambda(k) = \operatorname{sign}\!\left[y(k)-r\right]\,/\,T + A\,\lambda(k+1) \label{eq:adjoint}\\
\frac{\partial L}{\partial u(k)} &= B\,\lambda(k+1) \label{eq:grad_u}
\end{align}

Este gradiente se retropropaga luego a través de la red (backprop estándar) y los pesos se actualizan con Adam con cosine LR annealing.

Intuición. Cada acción $u(k)$ no afecta solo al error de ese instante: arrastra toda la trayectoria posterior de $y$. Derivar muestra a muestra, por tanto, no basta. El co-estado $\lambda(k)$ resume justo eso, cuánto pesa el estado del instante $k$ en el error acumulado hasta el final del episodio. El cálculo va en tres pasos:
\begin{enumerate}
  \item Pasada hacia adelante: se simula el episodio completo con la red actual y se guardan $y(k)$, $u(k)$ y las activaciones de la red.
  \item Pasada hacia atrás: se parte de $\lambda(T)$ (el error en el instante final) y se recorre el tiempo en sentido inverso; en cada paso se suma el error local más la contribución que llega del futuro, $A\,\lambda(k+1)$. Es la misma idea de \emph{backpropagation through time}.
  \item Gradiente: $\partial L/\partial u(k)=B\,\lambda(k+1)$ indica cómo mover cada acción para reducir el IAE; ese valor se inyecta como error en la capa de salida de la red y se retropropaga con backprop estándar.
\end{enumerate}
En la práctica equivale a derivar la simulación completa respecto a las acciones de la red. No hace falta un controlador que le indique la ``respuesta correcta'': la red aprende directamente de cuánto error produce sobre la planta.

\subsection*{Fase~1: Servoregulador}

Estado de entrada: $[e,\;\Delta e,\;\sigma,\;d]$ con $\sigma(k)=\sum_{j=0}^{k}e(j)$ (integral acumulada). La integral distingue error estacionario acumulado de estado puro; sin ella la red no puede eliminar el offset. Arquitectura: $4\!\to\!64\!\to\!32\!\to\!16\!\to\!8\!\to\!1$ (sigmoide en salida). Entrenamiento: 800~épocas × $B=32$ escenarios aleatorios ($y_0\sim U(0,0.4)$, perturbación aleatoria). Pérdida: $L=\text{IAE}+\lambda_u\sum|\Delta u|^2$, con $\lambda_u=0.01$ para penalizar cambios bruscos del mando.

\subsection*{Fase~2: Coordinación aware}

El mismo principio adjunto se extiende a la planta R2ET de dos esteras. Estado: $[e_1,\,e_2,\,\Delta e_1,\,\Delta e_2,\,\sigma_1,\,\sigma_2,\,b,\,d]$ (8 entradas). Arquitectura más profunda que en el servo, $8\!\to\!96\!\to\!64\!\to\!32\!\to\!16\!\to\!8\!\to\!3$ (tres salidas sigmoide: $u_1$, $u_2$, $r_1$), por la mayor complejidad del problema de dos esteras. El co-estado se calcula para cada estera por separado ($\lambda_1$, $\lambda_2$) y el gradiente respecto a $r_1$ incorpora el efecto del reparto sobre la alimentación de cada estera. Entrenamiento: 1200~épocas × $B=16$ escenarios ($n_0\sim U(0.4,2.6)$, atasco aleatorio con probabilidad $0.55$, demanda $\pm15\%$).

\artifactfigure{figures/plots/fig_dl_training.png}{0.90}{Curvas de entrenamiento. Izq.: Fase~1 servo: IAE normalizado por optimización directa (método adjunto, 800~épocas × 32~escenarios). Der.: Fase~2 aware: IAE normalizado por método adjunto de dos esteras (1200~épocas × 16~escenarios). Escala logarítmica.}

% Tabla de resultados - Act08 R2ET
\begin{table}[H]
\centering
\caption{Entrenamiento Deep Learning: optimización directa sobre planta (método adjunto).}
\label{tab:mlp_training}
\small
\resizebox{\linewidth}{!}{%
\begin{tabular}{@{}lrrrrrrrr@{}}
\toprule
Fase & Arquitectura & Épocas & Escenarios & Val. & Test & IAE entren. & IAE val. & IAE test \\
\midrule
Servo & 4-64-32-16-8-1 & 800 & 25600 & 50 & 50 & 0.00406 & 0.00432 & 0.00432 \\
Aware & 8-96-64-32-16-8-3 & 1200 & 19200 & 30 & 30 & 0.205 & 0.173 & 0.173 \\
\bottomrule
\end{tabular}%
}
\end{table}


\subsection*{Métricas de aprendizaje aplicables al control con redes profundas}

La red DL produce salidas continuas $u\in[0.05,\,1.0]$ (velocidad de actuador), por lo que actúa como regresor, no como clasificador. Las métricas propias de clasificación (AUC-ROC, precisión, sensibilidad, F1-score, matrices de confusión) no aplican: requieren etiquetas binarias y probabilidades de clase, que no existen aquí. Los indicadores correctos para un regresor de control son:

\begin{longtable}{@{}P{0.28\linewidth}P{0.30\linewidth}P{0.30\linewidth}@{}}
\caption{Métricas de aprendizaje aplicadas a la red de control.}
\label{tab:ml_metrics}\\
\toprule
Métrica & Servo ($4{\to}1$) & Aware ($8{\to}3$) \\ \midrule \endfirsthead
\toprule Métrica & Servo & Aware \\ \midrule \endhead
$L_{train}$ (IAE normaliz.) & 0.00406 & 0.205 \\
$L_{val}$ (IAE normaliz.)   & 0.00432 & 0.173 \\
Brecha de generalización $\Delta_{gen}$ & $\mathbf{6.4\%}$ & $-15.6\%$ \\
Veredicto sobreajuste & Sin sobreajuste & Sin sobreajuste ($L_{val}<L_{train}$) \\
Épocas hasta convergencia & $\approx 600$ de 800 & $\approx 1000$ de 1200 \\
Inicialización & He (\texttt{sqrt(2/fan\_in)}) & He (\texttt{sqrt(2/fan\_in)}) \\
Optimizador & Adam, LR cosine & Adam, LR cosine \\
\bottomrule
\end{longtable}

La brecha de generalización $\Delta_{gen}=\left(L_{val}-L_{train}\right)/L_{train}$ mide si la red memoriza los escenarios de entrenamiento en lugar de aprender la política general. Un valor $<5\%$ indica generalización excelente; entre $5\%$-$20\%$ es aceptable; $>30\%$ indica sobreajuste severo. Para la Fase~1, $\Delta_{gen}=6.4\%$: la red generaliza bien a escenarios no vistos (dentro del rango aceptable). Para la Fase~2, $\Delta_{gen}$ es \emph{negativo} ($L_{val}<L_{train}$): no hay sobreajuste, y la red profunda con 1200~épocas generaliza bien al conjunto de validación.

\section{Q-Learning tabular}

Actualización TD(0) con política $\varepsilon$-greedy \cite[cap.~6]{sutton2018reinforcement}:
\[
Q(s,a) \leftarrow Q(s,a) + \alpha\!\left[r + \gamma\max_{a'}Q(s',a') - Q(s,a)\right]
\]

Durante el entrenamiento, $\varepsilon$ decae linealmente de $\varepsilon_0$ a $\varepsilon_{min}$; en despliegue $\varepsilon=0$ (política greedy pura).

El espacio de estados se discretiza de forma no uniforme con mayor resolución cerca del objetivo:
\begin{itemize}
  \item Servo: $s=(e_{bin},\,\Delta e_{bin},\,d_{bin})$ → $10\times5\times2=100$ estados; 8~acciones de velocidad discreta.
  \item Aware: $s=(n_{1,bin},\,n_{2,bin},\,d_{bin})$ → $9\times9\times2=162$ estados; 11~combinaciones de $(u,\,r_1,\,a)$.
\end{itemize}

La recompensa es negativa (un coste que el agente aprende a minimizar) y suma varios términos, cada uno con un peso que fija su importancia relativa:
\[
r_{servo} = \underbrace{-|e|}_{\text{error}} \underbrace{-\,0.3\,e^2}_{\text{penaliza error grande}} \underbrace{-\,0.05\,|\Delta u|}_{\text{suaviza el mando}} \underbrace{-\,0.5\max(0,\,y-(r{+}0.08))^2}_{\text{castiga el sobreimpulso}}
\]
\[
r_{aware} = \underbrace{-1.5\,\text{IAE}}_{\text{seguir }n^*} \underbrace{-\,30\,\text{viol.}}_{\text{nunca desbordar}} \underbrace{-\,8\,\text{alarma}}_{\text{evitar zona crítica}} \underbrace{-\,0.20\,|n_1-n_2|}_{\text{equilibrar esteras}}
\]
Los pesos altos de \emph{viol.} y \emph{alarma} (30 y 8) hacen que el agente aprenda, ante todo, a no superar la capacidad; el resto afina el seguimiento y el equilibrio.

\artifactfigure{figures/plots/fig_ql_training.png}{0.90}{Recompensa media por episodio (suavizada) y decaimiento de $\varepsilon$ durante el entrenamiento Q-Learning: servo (izq.) y aware (der.).}

% Tabla de resultados - Act08 R2ET
\begin{table}[H]
\centering
\caption{Parámetros y resultados del entrenamiento Q-Learning.}
\label{tab:ql_training}
\small
\resizebox{\linewidth}{!}{%
\begin{tabular}{@{}lrrrrrrrrr@{}}
\toprule
Fase & Estados & Acciones & Episodios & Horizonte & $\alpha$ & $\gamma$ & $\epsilon_0$ & $\epsilon_{min}$ & $r_{final}$ \\
\midrule
Servo & 100 & 8 & 600 & 100 & 0.15 & 0.95 & 0.6 & 0.02 & -0.038 \\
Aware & 162 & 11 & 3000 & 160 & 0.18 & 0.92 & 0.65 & 0.02 & -0.28 \\
\bottomrule
\end{tabular}%
}
\end{table}


\subsection*{Métricas de aprendizaje aplicables al Q-Learning tabular}

El agente Q-Learning no produce probabilidades de clase: selecciona una acción discreta mediante $\pi(s)=\arg\max_a Q(s,a)$. Las métricas de clasificación (ROC, AUC, FP/FN) no aplican. Los indicadores correctos para un agente de RL tabular son:

\begin{longtable}{@{}P{0.40\linewidth}P{0.24\linewidth}P{0.24\linewidth}@{}}
\caption{Métricas de convergencia del Q-Learning tabular.}
\label{tab:ql_metrics}\\
\toprule
Métrica & Servo & Aware \\ \midrule \endfirsthead
\toprule Métrica & Servo & Aware \\ \midrule \endhead
Recompensa media últimos 30~ep. & $-0.038$ & $-0.280$ \\
$\varepsilon$ final (exploración residual) & 0.02 & 0.02 \\
Episodios hasta convergencia estable & $\approx400$ de 600 & $\approx2000$ de 3000 \\
IAE de despliegue (greedy, $\varepsilon=0$) & 1.80 & 25.75 \\
Fracción de estados visitados ($\geq1$ vez) & $\approx85\%$ & $\approx78\%$ \\
\bottomrule
\end{longtable}

Convergencia. La recompensa media de $-0.038$ para el servo indica baja penalización acumulada: el agente mantiene el error pequeño en la mayoría de las transiciones. El valor $-0.280$ del aware refleja el mayor espacio de estados y la penalización multi-componente (IAE + violaciones + alarma + desequilibrio). Ambos agentes alcanzan $\varepsilon=\varepsilon_{min}=0.02$ al final del entrenamiento, lo que significa que el $98\%$ de las decisiones son greedy.

Sensibilidad a hiperparámetros. El parámetro crítico del Q-Learning es el factor de descuento $\gamma$: valores altos ($\gamma\to1$) priorizan recompensas futuras, necesario para un sistema de control donde el objetivo es el comportamiento sostenido. Con $\gamma=0.95$ (servo) y $\gamma=0.92$ (aware), la política da más peso a los próximos 20 y 12~pasos respectivamente, que corresponden aproximadamente al tiempo de establecimiento del sistema. La tasa $\alpha$ controla la velocidad de actualización: $\alpha=0.15$--$0.18$ es un compromiso entre velocidad de convergencia y estabilidad de la tabla Q.

Comparación con métricas de clasificación. Si se quisiera interpretar la política QL como clasificador binario (``¿evita alarma sí/no?''), la tasa de éxito sería del $100\%$ (cero alarmas en el despliegue), lo que daría AUC=1.0 y sensibilidad=1.0. Sin embargo, este cálculo no aporta información útil porque la capacidad del supervisor garantiza ese resultado con independencia del controlador. La métrica discriminante real es el IAE de despliegue.
